\section{Introduction}
\label{sec:intro}

Ask any engineer at an oil refinery or chemical plant how their facility
is documented and the answer will invariably involve P\&IDs ---
Piping and Instrumentation Diagrams that encode every valve, pump,
instrument, and pipe connection in the plant.
Keeping an accurate, queryable digital model of this information is
increasingly important for regulatory compliance, predictive maintenance,
and the growing industry interest in digital
twins~\cite{toghraei2019}.
The difficulty is that most P\&IDs exist only as raster images, and
converting them into structured graphs by hand is both expensive and slow.

Automating this conversion is an active research problem, but it has run
into a hard constraint: data.
Plant operators cannot share engineering drawings externally, as these
documents contain sensitive process information, and the entire field has
consequently been working from a single public benchmark of just 12
annotated real-world images~\cite{sturmer2025}.
Stürmer~\etal~\cite{sturmer2025} recently showed that an end-to-end
transformer substantially outperforms classical modular pipelines on this
benchmark, but their model requires 60 real P\&IDs for training.
They also note in passing that training on synthetic data alone
``suffers significantly'' --- without elaborating on why.

That observation deserves closer inspection.
Generating synthetic training data is the natural fallback when
annotated real data is scarce, and it has worked well in robotics,
medical imaging, and remote sensing.
Why does it fail here?
Existing synthetic generators~\cite{paliwal2021,sturmer2025} work by
randomly placing symbol templates on a blank canvas and connecting them
with simple lines.
The output looks superficially like a P\&ID, but the underlying graph
topology is unrealistic: symbols cluster uniformly, junction degrees
follow a narrow range, and the spatial structure of a real process flow
is absent.
A model trained on such data encounters real P\&IDs as a genuinely
foreign distribution, not merely a different visual style.

We take a different approach.
Rather than generating from scratch, we use real P\&IDs as structural
seeds.
Our generation pipeline extracts the physical component graph from each
seed, applies spatial and symbolic perturbations that respect topological
constraints, and renders the result into a new annotated image.
Coupled with Weisfeiler-Lehman and perceptual hashing to filter
structurally or visually duplicate outputs, this process yields \ours{}:
665 synthetic P\&IDs whose graph-level statistics match those of real
drawings while remaining visually distinct from their seeds.

Combining \ours{} with a patch-based adaptation of
Relationformer~\cite{shit2022} --- patches are needed because P\&ID
symbols are too small at full-image scale for reliable detection ---
we achieve \textbf{63.8\% edge mAP} on OPEN100 without any real P\&IDs
in the training set.
This is 30 percentage points better than training on template-generated
synthetic data, 18 points above the modular baseline that does use real
data, and within 8 points of training on real P\&IDs directly.
Adding a small number of real images to the training mix (E3) closes the
remaining gap to within 2.8 pp of published state of the art.

Our contributions are as follows.
We release \ours{} as the first publicly available topology-preserving
synthetic P\&ID dataset with full graph-level annotation, together with
the generation pipeline and trained models.
We provide controlled experimental evidence that generation quality,
rather than model architecture, is the primary driver of
synthetic-to-real performance.
Finally, we establish a rigorous evaluation protocol on OPEN100 ---
5-fold cross-validation with 3 independent random seeds --- as a
reproducible benchmark for future work in this data-scarce setting.
